\documentclass[12pt,a4paper]{article}

\usepackage{amsmath, amsthm, amssymb, amsfonts}
\usepackage{mathrsfs}
\usepackage{hyperref}
\usepackage{geometry}
\usepackage{enumitem}
\usepackage{xcolor}

\geometry{margin=2.5cm}

\newtheorem{theorem}{Theorem}[section]
\newtheorem{proposition}[theorem]{Proposition}
\newtheorem{lemma}[theorem]{Lemma}
\newtheorem{corollary}[theorem]{Corollary}
\newtheorem{definition}[theorem]{Definition}
\newtheorem{remark}[theorem]{Remark}
\newtheorem{assumption}[theorem]{Assumption}
\newtheorem{openproblem}[theorem]{Open Problem}

\newcommand{\Ai}{\operatorname{Ai}}
\newcommand{\lam}[1]{\lambda_{#1}^{(c)}}
\newcommand{\pswf}[1]{\psi_{#1}^{(c)}}
\newcommand{\norm}[1]{\left\|#1\right\|}
\newcommand{\FAi}{F_{\mathrm{Ai}}}

\title{%
  \textbf{B-strong Reduction Lemma}\\[0.4em]
  \large Airy--Galerkin Reduction, Regime Analysis,\\
  and Failure-Mode Characterisation\\
  \normalsize\textit{(Revised: Regime Misalignment Correction, May 2026)}
}
\author{\textit{Working note -- prolate-gram-coercivity programme}}
\date{May 2026}

\begin{document}
\maketitle

\begin{abstract}
We carry out a systematic reduction of Assumption~B-strong
($P_{kl} \leq C_2 c^{1/2}$, Paper~VI) to a question about
Airy-tail integrals on the Galerkin discretisation scale.

\medskip\noindent
\textbf{Revision note.}
A numerical calibration (Section~\ref{sec:misalignment}) revealed
a structural misalignment in the initial regime decomposition:
the analytically identified ``transition regime''
$x \in [A_0, (\log c)^{2/3}]$ is disjoint from the
spectral support of the Galerkin truncation $k \leq N = \lfloor Ac^{1/3}\rfloor$.
The Galerkin indices are entirely confined to the
\emph{superconcentration regime} $x_k \ll -1$,
where a fundamentally different asymptotic applies.
The regime analysis and open problem have been
recalibrated accordingly.

\medskip\noindent
This note documents a \emph{reduction result and a domain
correction}, not a proof of B-strong.
Its contribution is the exact localisation of the
spectral domain where B-strong is non-trivial,
and the structural reason why Airy-WKB methods calibrated
to the spectral edge cannot reach it.
\end{abstract}

\tableofcontents
\bigskip\hrule\bigskip

%% ============================================================
\section{Setup and Notation}
%% ============================================================

Throughout, $c > 0$ is the time-bandwidth product,
$N = \lfloor Ac^{1/3} \rfloor$ the Galerkin truncation,
and $\pswf{k}$, $\lam{k}$ the $k$-th PSWF and Slepian
concentration eigenvalue.
The Airy-scaled index is
\[
  x_l = (l - N_{\mathrm{Sh}})(c/2)^{-1/3},
  \qquad N_{\mathrm{Sh}} = \frac{2c}{\pi}.
\]

\paragraph{Two fundamental scales.}
The analysis involves two scales that are often conflated
but are asymptotically disjoint for $B$-strong:
\begin{itemize}[nosep]
  \item \textbf{Spectral-edge scale:} $k \sim N_{\mathrm{Sh}} = 2c/\pi$,
        giving $x_k \sim O(1)$.  This is where the Airy approximation
        to the eigenvalue distribution is valid, and where Regimes~I--III
        below are relevant.
  \item \textbf{Galerkin scale:} $k \leq N = \lfloor Ac^{1/3}\rfloor
        \ll N_{\mathrm{Sh}}$, giving $x_k = (k - N_{\mathrm{Sh}})(c/2)^{-1/3}
        \approx -N_{\mathrm{Sh}} \cdot (c/2)^{-1/3} \sim -c^{2/3}$.
        This is the \emph{superconcentration regime}, far to the left
        of the Airy window.
\end{itemize}
Numerically: for $c = 100$, $A = 2$, one has $N = 9$,
$N_{\mathrm{Sh}} \approx 63.7$, and $x_k \in [-17, -15]$
for $k \in [0, N]$.  The transition regime $x \sim (\log c)^{2/3}
\approx 2.8$ corresponds to $k \approx 74$, far outside the
Galerkin range.

\medskip
The Airy tail integral, its derivative, and the key ratio:
\[
  S(x) := \int_x^\infty \Ai(t)^2\,dt,
  \quad S'(x) = -\Ai(x)^2,
  \quad \Phi(x) := \frac{\Ai(x)^2}{S(x)^2}
  = \frac{d}{dx}\!\left(\frac{1}{S(x)}\right).
\]
The prefactor to be bounded (Assumption B-strong, Paper~VI):
\[
  P_{kl} := c \cdot \frac{|\lam{l} - \lam{k}|}{(1-\lam{k})(1-\lam{l})}.
\]
Throughout $d := |k-l|$, and we write $P_{k,k+d}$ for $P_{kl}$.

\medskip\noindent
We use the following results from \texttt{airy\_discrete\_stability\_lemma.tex}
without reproof:
\begin{itemize}[nosep]
  \item \textbf{Proposition~U}: $|\lam{l} - \FAi(x_l)| \leq C_1 c^{-2/3}$
        uniformly in $l \leq N$.
  \item \textbf{Proposition~U$'$}:
        $|(\lam{l}-\lam{l+1}) - (\FAi(x_l)-\FAi(x_{l+1}))|
        \leq C_2 c^{-2/3}$.
  \item \textbf{Airy spacing}: $x_{l+1} - x_l = (c/2)^{-1/3}(1+O(c^{-1}))$.
\end{itemize}

%% ============================================================
\section{Regime Misalignment Correction}
\label{sec:misalignment}
%% ============================================================

\begin{remark}[Regime Misalignment -- Structural Revision]
\label{rem:misalignment}
The initial version of this note analysed B-strong via a
three-regime decomposition calibrated to the \emph{spectral edge}
($x_k \in [-A_0, (\log c)^{2/3}]$).
Numerical verification revealed that this decomposition is
\emph{spectrally irrelevant} for the Galerkin indices $k \leq N$:

\begin{enumerate}
  \item \textbf{Scale separation.}
        The Galerkin truncation gives $k \leq N = \lfloor Ac^{1/3}\rfloor$,
        while the spectral transition region lies at
        $k \sim N_{\mathrm{Sh}} = 2c/\pi$.  Since $c^{1/3} \ll c/\pi$,
        these scales are asymptotically disjoint:
        \[
          \frac{N}{N_{\mathrm{Sh}}} = \frac{\lfloor Ac^{1/3}\rfloor}{2c/\pi}
          \sim \frac{A\pi}{2}\,c^{-2/3} \to 0.
        \]

  \item \textbf{Actual $x_k$ for Galerkin indices.}
        For $k \leq N$:
        \[
          x_k = (k - N_{\mathrm{Sh}})(c/2)^{-1/3}
          \approx -N_{\mathrm{Sh}}(c/2)^{-1/3}
          = -\frac{2c}{\pi}\cdot\frac{2^{1/3}}{c^{1/3}}
          = -\frac{2^{4/3}}{\pi}\,c^{2/3} \ll -1.
        \]
        The Galerkin indices live at $x_k \sim -c^{2/3}$,
        far to the left of the Airy window $|x| \lesssim 1$
        and of the transition regime $x \sim (\log c)^{2/3}$.

  \item \textbf{Consequence for the regime decomposition.}
        Regimes~I (Airy window) and~III (transition) were
        analysed at $x \in [-A_0, (\log c)^{2/3}]$.
        Since no Galerkin index has $x_k$ in this range,
        these regimes, while analytically correct,
        are \emph{spectrally empty} for B-strong.
        The divergence $c^{7/6}$ found in Regime~III is
        not a counterexample to B-strong: it arises from
        evaluating the formula $P_{kl}$ at indices $k$
        that do not appear in the Galerkin sum.

  \item \textbf{What was actually shown.}
        The analytic failure -- the impossibility of closing
        the $\Phi$-bound by pointwise methods at
        $x \sim (\log c)^{2/3}$ -- is a valid structural result
        about the spectral-edge regime.  It implies that any
        attempt to extend the Galerkin analysis to
        $k \sim N_{\mathrm{Sh}}$ via these methods would fail.
        For the actual Galerkin indices, a different
        asymptotic is operative.
\end{enumerate}
\end{remark}

%% ============================================================
\section{The Superconcentration Regime}
\label{sec:superconc}
%% ============================================================

For $k \leq N = \lfloor Ac^{1/3}\rfloor$, the Airy-scaled index
satisfies $x_k \sim -c^{2/3}$.
In this regime the PSWF eigenvalues satisfy
\[
  \lam{k} = 1 - \varepsilon_k,
  \qquad \varepsilon_k \ll 1,
\]
where $\varepsilon_k$ is exponentially small in $c$.
This is the \emph{superconcentration regime}:
all Galerkin eigenvalues are exponentially close to~$1$.

\begin{proposition}[B-strong in the superconcentration regime]
\label{prop:superconc}
For $k, l \leq N$ with $d = |k-l|$ and $\varepsilon_k, \varepsilon_l
\ll 1$:
\[
  P_{kl}
  = c \cdot \frac{|\varepsilon_l - \varepsilon_k|}{\varepsilon_k \varepsilon_l}.
\]
B-strong ($P_{kl} \leq Cc^{1/2}$) is therefore equivalent to:
\[
  \frac{|\varepsilon_l - \varepsilon_k|}{\varepsilon_k \varepsilon_l}
  \leq C c^{-1/2}.
\]
This is a statement about the \emph{relative variation of
exponentially small eigenvalue gaps} in the superconcentration
regime, not about Airy-WKB interpolation.
\end{proposition}

\begin{remark}[Recalibrated open problem]
\label{rem:recalibrated}
The correct formulation of the open problem for B-strong is:

\medskip
\textit{Show that for all $k, l \leq N = \lfloor Ac^{1/3}\rfloor$
and $d = |k-l| \leq c^{1/6}$, the exponentially small
residuals $\varepsilon_k = 1 - \lam{k}$ satisfy}
\[
  \frac{|\varepsilon_{k+d} - \varepsilon_k|}{\varepsilon_k \,
  \varepsilon_{k+d}} \leq C c^{-1/2}.
\]

\noindent
The relevant asymptotic tool is not the Airy CDF $S(x)$ at
the spectral edge, but rather the asymptotics of
$\varepsilon_k = 1 - \lam{k}$ in the regime $k \ll N_{\mathrm{Sh}}$,
where $\lam{k}$ is superexponentially concentrated near $1$.
The appropriate framework is therefore either:
\begin{itemize}[nosep]
  \item \textbf{Eigenvalue exponential asymptotics}
        (Landau--Pollak, Slepian): precise estimates for
        $1 - \lam{k}$ as $k \to 0$ or $k \ll N_{\mathrm{Sh}}$;
  \item \textbf{Resolvent/Green-function methods}
        for the prolate operator restricted to the
        superconcentration sector;
  \item \textbf{Operator coercivity directly}
        on the Galerkin subspace spanned by $\pswf{0}, \ldots, \pswf{N}$,
        bypassing the eigenvalue-ratio representation entirely.
\end{itemize}
\end{remark}

%% ============================================================
\section{Retention: What the Airy Analysis Correctly Shows}
\label{sec:retention}
%% ============================================================

Despite the regime misalignment, the Airy-WKB analysis of
Sections~2--5 (original version) establishes valid results
in their own domain:

\begin{enumerate}
  \item \textbf{Telescope identity (Proposition~\ref{prop:telescope_retained}).}
        Valid for all $k, l$ where $1 - \lam{k} \approx S(x_k)$,
        i.e.\ for $k$ near the spectral edge $k \sim N_{\mathrm{Sh}}$.
        It gives a sharp reduction of $P_{kl}$ to a
        difference of $1/S$, with controlled error.

  \item \textbf{Failure-Mode characterisation.}
        The impossibility of controlling $\Phi(x)$ uniformly
        on $[A_0, (\log c)^{2/3}]$ is a genuine obstruction
        for any B-strong-type bound that might be formulated
        for indices $k \sim N_{\mathrm{Sh}}$.
        It shows that the spectral-edge problem would require
        correlated Airy-kernel methods beyond pointwise WKB.

  \item \textbf{Structural principle.}
        The identity $d(1/S)/dx = \Phi(x)$ and the
        Wronskian relation $S' = -\Ai^2$ characterise the
        correct algebraic invariant for any future
        spectral-edge B-strong analysis.
\end{enumerate}

\begin{proposition}[Telescope identity -- retained]
\label{prop:telescope_retained}
For $k, l$ with $|x_k|, |x_l| \leq C_0$ (spectral-edge regime):
\[
  P_{k,k+d}
  = c^{2/3}\!\left(\frac{1}{S(x_{k+d})} - \frac{1}{S(x_k)}\right)
  + \mathcal{E}_{k,d},
  \qquad |\mathcal{E}_{k,d}| \leq C_{\mathcal{E}}\, c^{1/3}.
\]
\end{proposition}

%% ============================================================
\section{Summary: Revised Logical Map}
%% ============================================================

\begin{center}
\renewcommand{\arraystretch}{1.5}
\begin{tabular}{p{3.2cm}p{5cm}p{2.5cm}p{2cm}}
\hline
\textbf{Item} & \textbf{Content} & \textbf{Domain} & \textbf{Status} \\
\hline
Telescope identity & $P_{kl} = c^{2/3}\Delta(1/S)+\mathcal{E}$ &
  spectral edge $x_k \sim 0$ & \checkmark\ exact \\
Failure-Mode (b) & $\Phi$ not uniform; Airy correlation required &
  spectral edge & \checkmark\ proven \\
Regime~I (Airy) & controlled via $\delta \leq c^{-1/6}$ &
  spectral edge & \checkmark\ closed \\
Regime~II (exp.) & controlled via direct $\varepsilon$ bound &
  spectral edge & \checkmark\ closed \\
Regime~III (transition) & $\Phi$ blows up; pointwise methods fail &
  spectral edge & proven obstruction \\
\hline
Galerkin B-strong & $\varepsilon_k^{-1}|\varepsilon_{k+d}-\varepsilon_k|
  \cdot \varepsilon_{k+d}^{-1} \leq Cc^{-1/2}$ &
  $x_k \sim -c^{2/3}$ & \textbf{open} \\
Superconc.\ asymptotics & Precise $\varepsilon_k = 1-\lam{k}$ for
  $k \ll N_{\mathrm{Sh}}$ & $k \leq N$ & \textbf{open} \\
Operator coercivity & Direct bound bypassing $P_{kl}$ &
  Galerkin subspace & \textbf{open} \\
\hline
\end{tabular}
\end{center}

\medskip\noindent
\textbf{Core revised diagnosis.}
B-strong for the actual Galerkin indices $k \leq N = O(c^{1/3})$
is a statement about \emph{exponentially small eigenvalue
residuals in the superconcentration regime},
not about Airy-WKB interpolation at the spectral edge.
The two problems are connected (both concern the prolate operator)
but require different asymptotic tools.
The Airy analysis of this note is valid and non-trivial in
its own domain (spectral edge), and correctly characterises
the failure of pointwise methods there.
For the Galerkin problem proper, the appropriate tools are
eigenvalue exponential asymptotics or direct operator
coercivity arguments on the concentrated subspace.

%% ============================================================
\section*{Relation to the Programme}
%% ============================================================

This note is a companion to Paper~VI (\texttt{paper6.tex}).
The regime misalignment identified here sharpens the
programme in two ways:
\begin{enumerate}[nosep]
  \item The Airy-WKB approach is the correct tool for any
        future analysis of B-strong-type bounds near
        $k \sim N_{\mathrm{Sh}}$ (spectral edge), including
        potential extensions of Paper~VI to that regime.
  \item For the Galerkin B-strong as stated (with $k \leq N
        = O(c^{1/3})$), the analysis must be redirected to
        superconcentration asymptotics.
        The most direct route is likely a
        direct coercivity estimate on the restricted
        operator $T_c|_{V_N}$, using the known
        exponential concentration of $\lam{k}$ near~$1$
        for $k \leq N$.
\end{enumerate}

%% ============================================================
\begin{thebibliography}{9}

\bibitem{paper6}
  \textit{Paper~VI: Galerkin Norm Estimates for the Prolate Phase Operator},
  \texttt{paper6.tex}, prolate-gram-coercivity programme, May~2026.

\bibitem{airy_lemma}
  \textit{Airy Discrete Stability Lemma},
  \texttt{airy\_discrete\_stability\_lemma.tex},
  prolate-gram-coercivity programme, 2026.

\bibitem{Slepian1978}
  D.~Slepian,
  ``Prolate Spheroidal Wave Functions, Fourier Analysis, and
  Uncertainty~V: The Discrete Case,''
  \textit{Bell Syst.\ Tech.\ J.}, 57 (1978), 1371--1430.

\bibitem{OsipovRokhlinXiao2013}
  A.~Osipov, V.~Rokhlin, H.~Xiao,
  \textit{Prolate Spheroidal Wave Functions of Order Zero},
  Springer, 2013.

\end{thebibliography}

\end{document}
